\section{有限秩线性映射}

设 A 是一个环，E 与 F 是 A-模。我们用 \(\operatorname{Hom}_{\mathrm{A}}^{\mathrm{f}}(\mathrm{E},\mathrm{F})\) 表示从 E 到 F 的像包含在 F 的一个有限生成子模中的线性映射所成的集合。它是 \(\operatorname{Hom}_{\mathrm{A}}(\mathrm{E},\mathrm{F})\) 的一个子群。当 A 是一个体时，它是从 E 到 F 的有限秩线性映射所成的集合（II, §7, No.~4, 第 298 页，定义 3）。置 \(\operatorname{End}_{\mathrm{A}}^{\mathrm{f}}(\mathrm{E}) = \operatorname{Hom}_{\mathrm{A}}^{\mathrm{f}}(\mathrm{E},\mathrm{E})\)。

设 E、F、G 是 A-模，\(u : E \to F\) 与 \(v : F \to G\) 是 A-线性映射。若 \(u \in \operatorname{Hom}_{\mathrm{A}}^{\mathrm{f}}(\mathrm{E}, \mathrm{F})\) 或 \(v \in \operatorname{Hom}_{\mathrm{A}}^{\mathrm{f}}(\mathrm{F}, \mathrm{G})\)，则 \(v \circ u\) 属于 \(\operatorname{Hom}_{\mathrm{A}}^{\mathrm{f}}(\mathrm{E}, \mathrm{G})\)。

用 \(\theta\) 表示从 \(\mathrm{E}^* \otimes_{\mathrm{A}} \mathrm{F}\) 到 \(\mathrm{Hom}_{\mathrm{A}}(\mathrm{E}, \mathrm{F})\) 的典范群同态（II, §4, No.~2, 第 271 页）；它把元素 \(x^* \otimes y\)（属于 \(\mathrm{E}^* \otimes_{\mathrm{A}} \mathrm{F}\)）映为映射 \(x \mapsto \langle x, x^* \rangle y\)（从 \(\mathrm{E}\) 到 \(\mathrm{F}\)）。

引理 1. --- 假设 A-模 F 是投射的。群同态 \(\theta\) 是单射，且它的像是 \(\mathrm{Hom}_{\mathrm{A}}^{\mathrm{f}}(\mathrm{E},\mathrm{F})\)。

由同前所引的推论，同态 \(\theta\) 是单射。它的像包含在 \(\operatorname{Hom}_{\mathrm{A}}^{\mathrm{f}}(\mathrm{E},\mathrm{F})\) 中。设 \(u \in \operatorname{Hom}_{\mathrm{A}}^{\mathrm{f}}(\mathrm{E},\mathrm{F})\)；让我们证明 u 属于 \(\theta\) 的像。

首先，假设 A-模 F 是自由的。设 \((f_{i})_{i\in\mathrm{I}}\) 是 F 的一个基，\((f_{i}^{*})_{i\in\mathrm{I}}\) 是 \(F^{*}\) 的与 \((f_{i})_{i\in\mathrm{I}}\) 对偶的基。设 J 是 I 的一个有限子集，使得 u 的像包含在由诸 \(f_{j}\)（\(j\in J\)）生成的 F 的子模中。我们有

\[
u (x) = \sum_ {j \in \mathrm{J}} \langle u (x), f _ {j} ^ {*} \rangle f _ {j} = \sum_ {j \in \mathrm{J}} \langle x, ^ {t} u (f _ {j} ^ {*}) \rangle f _ {j}\tag{1}
\]

对每个 \(x\in \mathrm{E}\) 成立，从而

\[
u = \theta \left(\sum_ {j \in \mathrm{J}} ^ {t} u \left(f _ {j} ^ {*}\right) \otimes f _ {j}\right).\tag{2}
\]

在一般情形，存在一个自由 A-模 L 与同态 \(i: F \to L\) 和 \(p: L \to F\)，使得 \(p \circ i = 1_{F}\)。同态 \(i \circ u\) 属于 \(\operatorname{Hom}_{\mathrm{A}}^{\mathrm{f}}(\mathrm{E}, \mathrm{L})\)。由上述，存在一个有限集 J，且对每个 \(j \in J\)，存在元素 \(x_{j}^{*}\)（属于 \(E^{*}\)）与 \(y_{j}\)（属于 L），使得我们有

\[
i (u (x)) = \sum_ {j \in \mathrm{J}} \langle x, x _ {j} ^ {*} \rangle y _ {j}
\]

对每个 \(x \in \mathrm{E}\) 成立。于是我们有

\[
u (x) = p \left(i (u (x))\right) = \sum_ {j \in J} \langle x, x _ {j} ^ {*} \rangle p \left(y _ {j}\right)
\]

对每个 \(x \in \mathrm{E}\) 成立，从而 \(u = \theta\big(\sum_{j \in \mathrm{J}} x_j^* \otimes p(y_j)\big)\)。

